% !TEX TS-program = pdflatex
% !TEX encoding = UTF-8 Unicode

% This is a simple template for a LaTeX document using the "article" class.
% See "book", "report", "letter" for other types of document.

\documentclass[11pt]{article} % use larger type; default would be 10pt

\usepackage[utf8]{inputenc} % set input encoding (not needed with XeLaTeX)

%%% Examples of Article customizations
% These packages are optional, depending whether you want the features they provide.
% See the LaTeX Companion or other references for full information.

%%% PAGE DIMENSIONS
\usepackage{geometry} % to change the page dimensions
\geometry{a4paper} % or letterpaper (US) or a5paper or....
% \geometry{margins=2in} % for example, change the margins to 2 inches all round
% \geometry{landscape} % set up the page for landscape
%   read geometry.pdf for detailed page layout information

\usepackage{graphicx} % support the \includegraphics command and options
\usepackage{amsmath}
% \usepackage[parfill]{parskip} % Activate to begin paragraphs with an empty line rather than an indent

%%% PACKAGES
\usepackage{booktabs} % for much better looking tables
\usepackage{array} % for better arrays (eg matrices) in maths
\usepackage{paralist} % very flexible & customisable lists (eg. enumerate/itemize, etc.)
\usepackage{verbatim} % adds environment for commenting out blocks of text & for better verbatim
\usepackage{subfig} % make it possible to include more than one captioned figure/table in a single float
% These packages are all incorporated in the memoir class to one degree or another...

%%% HEADERS & FOOTERS
\usepackage{fancyhdr} % This should be set AFTER setting up the page geometry
\pagestyle{fancy} % options: empty , plain , fancy
\renewcommand{\headrulewidth}{0pt} % customise the layout...
\lhead{}\chead{}\rhead{}
\lfoot{}\cfoot{\thepage}\rfoot{}

%%% SECTION TITLE APPEARANCE
\usepackage{sectsty}
\allsectionsfont{\sffamily\mdseries\upshape} % (See the fntguide.pdf for font help)
% (This matches ConTeXt defaults)

%%% ToC (table of contents) APPEARANCE
\usepackage[nottoc,notlof,notlot]{tocbibind} % Put the bibliography in the ToC
\usepackage[titles,subfigure]{tocloft} % Alter the style of the Table of Contents
\renewcommand{\cftsecfont}{\rmfamily\mdseries\upshape}
\renewcommand{\cftsecpagefont}{\rmfamily\mdseries\upshape} % No bold!

%%% END Article customizations

%%% The "real" document content comes below...

\title{Thermal}
%\author{Giuseppe Romano}
%\date{} % Activate to display a given date or no date (if empty),
         % otherwise the current date is printed 
\begin{document}
\maketitle

\section{Getting started}

Thermal module computes the balance between the heat generated and the heat dissipated. Here we get the temperature profile of a rectangular device with an hotspot in the center.

The device block specifies the material (Silicon) and the physical regions. In this case we have the HotSpot and the Base region,

\begin{verbatim}
Device
{
   material = Si
  
   Region HotSpot{}
   
   Region Base{}  
}
\end{verbatim}
Let's now declare the thermal module. The temperature is fixed at the left and right side of the domain. The hotspot is included only in the HotSpot region.

\begin{verbatim}
Module thermal
{
  Physics
  {

    Contact left 
    {
       type = heat_reservoir
       temperature = 300
    }
    
    Contact right 
    {
       type = heat_reservoir
       temperature = 300
    }

   HeatSource constant
   {
      regions = HotSpot
       H = 1e10
   }
  }
}
\end{verbatim}
In the simulation region we write which modules should be solved. 
\begin{verbatim}
Simulation{solve = thermal}
\end{verbatim}
The temperature map is shown below
\begin{figure}[h!]
\begin{center}
\includegraphics[width=10.0 cm]{hotspot.png}
\label{default}
\end{center}
\end{figure}

\section{Example: Joule's heat in a 2D device}

In this example we will see how to perform a thermal-driftdiffusion self-consistent simulation.
We include the Seebeck and Peltier effects. The heat conduction through the environment is modelled by adding an air region all around the diode. We assume that far from the junction the system is in thermal equilibrium with a thermal bath at 300 K. The heat dissipation through the substrate is taken into acount by introducing a  thermal surface resistance. Furthermore, we rely on the cylindrical symmetry with respect to the axis growth and consider only a 2D slice of the system. In this way we simulate a 3D device by performing a 2D simulation (with much less computational time). The mesh corresponding to the whole system (diode plus air region) is shown below.
\begin{figure}[h!]
\begin{center}
\includegraphics[width=10.0 cm]{tut_09_mesh.png}
\label{default}
\end{center}
\end{figure}
Region n-side corresponds to the n-doped region situated at the bottom of the device, while the region p-side refers to the top p-doped region. Air is modelled with Region Air.
\begin{verbatim}
Device
{

 meshfile = tut_09.msh
 
 dimension = 2

  material = Si
  
  Region n_side { doping = 1e18 doping_type = donor }
 
  Region p_side{ doping = 1e19 doping_type = acceptor}

  Region air{}

}
\end{verbatim}
The drift diffusion simulation is performed only over the device domain (excluding the Air region). Therefore, in option subsection we have to indicate \emph{physical\_regions = (n\_side,p\_side)} . The mobility models used for both electrons and holes are doping dependent. 

\begin{verbatim}
electron_mobility doping_dependent{}

hole_mobility doping_dependent{}
\end{verbatim}
The connection with the thermal simulation is given by the keyword \emph{thermal\_simulaation = tt}. The thermoelectric power model used can be included with \emph{ thermoelecric diffusivity\_model} (see the theory section).  Direct bias applied to the diode is obtained by using the following boundary conditions

\begin{verbatim}

Contact anode{@Vb[0.0]}

Contact cathode{}

\end{verbatim}
Unlike the drift diffusion case, the thermal simulation has to be performed over the whole simulation domain (including the Air region). We can refer to all regions simply by writing 
\emph{physical\_regions = all}.
The connection with the drift diffusion simulation is specified with \begin{verbatim} HeatSource joule {DD_simulation_name = dd} \end{verbatim}
The Dirichlet boundary condition is set with 
\begin{verbatim}
Contact reservoir
{
type = heat_reservoir
temperature = 300
}
\end{verbatim}
whereas the surface thermal resistance is applied with

\begin{verbatim}
Contact dissipator
{
  r_surf = 0.5
  type = surface_resistance
  temperature = 300
}
\end{verbatim}
The selfconsistent simulation solves the thermal and the drift diffusion simulations (named tt and dd respectively) in a self consistent way.

\begin{verbatim}
Module Selfconsistent
{
 flavour = relaxation 
 simulations = (dd,tt)
 abs_tolerance = 1e-5
 monitor = true
}
\end{verbatim}
The sweep section runs the selfconsistent simulation for each anode bias step.

\begin{verbatim}
Module sweep
{
simulation = selfconsistent
variable = Vb
start = 0.0
stop = 1.0
steps = 10
}
\end{verbatim}


\begin{verbatim}
Simulation
{
temperature = 300
solve = sweep
symmetry = cylindrical
}
\end{verbatim}
The temperature map, the heat source and the thermoelectri power are shown below.
\begin{figure}[h!]
\begin{center}
\includegraphics[width=10.0 cm]{tut_09_temperature.png}
\label{default}
\end{center}
\end{figure}
\begin{figure}[h!]
\begin{center}
\includegraphics[width=10.0 cm]{tut_09_HEAT.png}
\label{default}
\end{center}
\end{figure}
\begin{figure}[h!]
\begin{center}
\includegraphics[width=10.0 cm]{tut_09_TEP.png}
\label{default}
\end{center}
\end{figure}

\section{Theory}

Details about irreversible thermodynamics can be found in [Wachutka, IEEE Trans CAD Cicui. Integr and Syst. (1990)].

The heat flux, within the diffusive approach, is given by the Fourier’s law, i. e.
\begin{equation}
J_i = - \kappa_{ij} \frac{\partial}{\partial x_j} T
\end{equation}
where $\kappa$ is thermal conductivity second-rank tensor (which is assumed temperature independent). The power balance is ensured by the continuity equation for the thermal flux
\begin{equation}
\frac{\partial}{\partial x_i} J_i = H
\end{equation}



\subsection{Boundary conditions}

The boundaries of the thermal simulation domain are set to ideally thermal insulator by default.
Dirichlet boundary conditions can be imposed with the keyword boundary \emph{heat\_researvoir}. 
\begin{verbatim}
   Contact reservoir
   {
     temperature = 300
   }
\end{verbatim}
A surface boundary resistance can be added with the keyword  \emph{surface\_resistance}. 
\begin{verbatim}
   Contact surface_resistance
   {
     r_surf = 0.01
     temperature = 300
   }
\end{verbatim}

\subsection{Heat sources}

The constant value heat source can be included with \emph{heat\_source constant}. 
Example:  
\begin{verbatim}
   HeatSource constant{H=1e6}. 
\end{verbatim}

The heat source must be in W/cm3. Heat generated by the Joule’s effect is included with \emph{heat\_source joule}. The name of the drift-diffusion simulation is given by \emph{dd\_simulation}. 
Example:  
\begin{verbatim}
   HeatSource joule{dd_simulation = dd}. 
\end{verbatim}

\newpage
\section{Reference Guide}


\begin{verbatim}
Heat Source/constant
Kind: bulk 
Name: HeatSource
Type: constant
database_section: none
Options: 
   H (double,0.0)

Heat Source/joule
Kind: bulk 
Name: HeatSource
Type: joule
database_section: none
Options: 
   dd_simulation (string,”driftdiffusion”)

Boundary/Heat reservoir
Kind: interface
Name: Contact
Type: heat_reservoir
database_section: none
Options: 
   temperature (positive double,300)

Boundary/Surface Thermal Resistance
Kind: interface
Name: Contact
Type: surface_resistance
database_section: none
Options: 
   r_surf (positive double 0.0)
   temperature (positive double (K), 300)

Boundary/Flux
Kind: interface
Name: Contact
Type: flux
database_section: none
Options: 
   Flux  (double vector (W/cm2, (0.0,0.0,0.0))


\end{verbatim}




\end{document}
